\documentclass[11pt]{article}
\usepackage[T1]{fontenc}
\usepackage{lmodern}
\usepackage[a4paper,margin=28mm]{geometry}
\usepackage{amsmath,amssymb,amsthm,booktabs,array,longtable}
\usepackage{microtype}
\usepackage[hidelinks]{hyperref}
\hypersetup{pdftitle={A small-component bound for odd primitive non-deficient numbers with five distinct prime factors},pdfauthor={Xiaojian Zeng}}
\newtheorem{theorem}{Theorem}
\newtheorem{lemma}[theorem]{Lemma}
\DeclareMathOperator{\rad}{rad}
\newcommand{\ph}{\varphi}
\title{A small-component bound for odd primitive\protect\\
non-deficient numbers with five distinct prime factors}
\author{Xiaojian Zeng\thanks{PBC School of Finance, Tsinghua University,
Beijing 100190, China. Email: \texttt{xiaojianzeng@foxmail.com}.}}
\date{}

\begin{document}
\maketitle
\begin{abstract}
Let $n=\prod_{i=1}^{5}p_i^{a_i}$ be odd and primitive non-deficient,
where the $p_i$ are distinct primes and every $a_i$ is positive.
We prove that $\min_i p_i^{a_i+1}<2\rad(n)$, establishing the
five-prime case of Zelinsky's small-component conjecture.
A quotient estimate reduces a hypothetical counterexample, uniformly in
the exponents, to the radical gaps $R-2\ph(R)=1$ or $3$.
A congruence excludes the first gap. The second is excluded by a complete
thirteen-row divisor certificate, whose six remaining ordered integral
pairs each contain a composite number. No bound on the primes or
exponents is assumed.
\end{abstract}

\section{Statement and context}

Write $\sigma(m)=\sum_{d\mid m}d$ and $h(m)=\sigma(m)/m$.
A positive integer is \emph{deficient} if $h(m)<2$ and
\emph{non-deficient} if $h(m)\ge2$. It is \emph{primitive
non-deficient} if it is non-deficient and every positive proper divisor
is deficient. This definition includes perfect numbers.
A component of $n$ means a full prime power $p^a\parallel n$.

Zelinsky conjectures that every primitive non-deficient number
$n=\prod_{i=1}^{k}p_i^{a_i}$ has an index $i$ satisfying
$p_i^{a_i+1}<2\rad(n)$ \cite[Conjecture~1]{Zelinsky2026}.
His Theorem~2 proves the bound $2k\rad(n)$, and his Lemma~4 establishes
the conjectured bound for odd numbers with three or four distinct prime
divisors. We establish the next case.

\begin{theorem}\label{thm:main}
Let $3\le p_1<p_2<p_3<p_4<p_5$ be primes, let
$a_1,\ldots,a_5$ be positive integers, and put
\[
 n=\prod_{i=1}^{5}p_i^{a_i},\qquad R=\prod_{i=1}^{5}p_i.
\]
If $\sigma(n)\ge2n$ and $\sigma(d)<2d$ for every positive proper
divisor $d\mid n$, then
\[
 \min_{1\le i\le5}p_i^{a_i+1}<2R.
\]
\end{theorem}

Thus the general bound $10R$ in this case can be replaced by $2R$.
The proof uses the deficiency of $n/3$ to reduce all exponent choices
to two equations for the square-free integer $R$. The substantive step
is the complete exclusion of those equations; a classification of all
eligible $n$ is unnecessary.

The broader setting has a long history. Dickson proved finiteness of
the odd primitive non-deficient numbers with a fixed number of distinct
prime divisors \cite{Dickson1913}. Cohen relates primitivity to the
excess of the divisor sum and the sizes of components
\cite[Theorem~1 and Remark~1]{Cohen1983}; his Lemma~1 also records
the strict increase of $h$ under proper divisibility.
Amato, Hasler, Melfi and Parton give algorithms enumerating primitive
abundant numbers with a fixed number $\Omega$ of prime factors counted
with multiplicity \cite{AmatoEtAl2019}.
Here $\Omega=\sum_i a_i$ is unbounded, whereas the radical has exactly
five prime factors. The product estimate used below already appears in
Zelinsky's proof of Lemma~9 \cite{Zelinsky2026}; neither that elementary
inequality nor the divisor-pair factorization is claimed as a separate
methodological advance.

\section{The quotient estimate}

Put
\[
 D=\prod_{i=1}^{5}(p_i-1)=\ph(R),\qquad
 E=R-2D,\qquad b_i=p_i^{a_i+1}.
\]
The prime-power divisor-sum formula gives
\begin{equation}\label{eq:abundancy}
 h(n)=\frac{R}{D}\prod_{i=1}^{5}\left(1-\frac1{b_i}\right)
 <\frac{R}{D}.
\end{equation}
Consequently $E$ is a positive odd integer. Since $t/(t-1)$ decreases
for $t>1$, the assumption $p_1\ge5$ would give
\[
 \frac{R}{D}\le
 \frac{5\cdot7\cdot11\cdot13\cdot17}{4\cdot6\cdot10\cdot12\cdot16}
 =\frac{85085}{46080}<2,
\]
contrary to \eqref{eq:abundancy} and $h(n)\ge2$. Thus $p_1=3$.

Suppose that $b_i\ge2R$ for all five indices. Directly computing the
prime-power factors of the proper divisor $n/3$ yields
\begin{equation}\label{eq:quotient}
 h(n/3)=\frac{R}{D}\left(1-\frac3{b_1}\right)
                   \prod_{i=2}^{5}\left(1-\frac1{b_i}\right).
\end{equation}
This formula remains valid when $a_1=1$: the local contribution is
$(3/2)(1-3/9)=1$, as required when the prime $3$ disappears.

For $0\le u_j\le1$ one has
\begin{equation}\label{eq:loss}
 \prod_j(1-u_j)\ge1-\sum_j u_j.
\end{equation}
Indeed, if $P$ is the product of the preceding factors, then
$0\le P\le1$ and $1-P(1-u)=1-P+Pu\le1-P+u$; induction applies.
The five losses in \eqref{eq:quotient} sum to at most $7/(2R)$.
Primitivity therefore implies
\begin{equation}\label{eq:gap}
 2>h(n/3)\ge\frac{R-7/2}{D}
             =2+\frac{E-7/2}{D}.
\end{equation}
It follows that $E<7/2$, so $E\in\{1,3\}$.

The value $E=1$ is impossible modulo $3$. If some $p_i$, $i\ge2$,
is $1$ modulo $3$, then $D\equiv E\equiv0\pmod3$.
Otherwise all four such primes are $2$ modulo $3$, whence
$D\equiv2\pmod3$ and $E\equiv-2D\equiv2\pmod3$.
The following lemma excludes the remaining value.

\section{Excluding the radical gap three}

\begin{lemma}\label{lem:radical}
No product $R$ of five distinct odd primes satisfies
$R-2\ph(R)=3$.
\end{lemma}

\begin{proof}
Such a product satisfies $R/\ph(R)>2$, so the preceding comparison
again forces $p_1=3$.
For an initial segment of the prime factors write
$A=\prod p$, $B=\prod(p-1)$ and $g=A-2B$.
Adjoining a prime $q$ changes the gap to
\begin{equation}\label{eq:prefix}
 Aq-2B(q-1)=qg+2B.
\end{equation}
If a nonempty proper initial segment has positive gap, then $g\ge1$
and $B\ge2$, so its next extension has gap greater than $3$.
Further extensions increase a positive gap. A zero gap is impossible
because $A$ is odd. Every proper initial segment must therefore have
$A<2B$.

If $p_2\ge15$, then
\[
 \frac{R}{\ph(R)}\le\frac32\left(\frac{15}{14}\right)^4<2,
 \qquad 3\cdot15^4=151875<153664=4\cdot14^4.
\]
Hence $p_2\in\{5,7,11,13\}$. For $p_2=p$ and $p_3\ge t$,
\begin{equation}\label{eq:cutoff}
 \frac{R}{\ph(R)}\le
 \frac{3p}{2(p-1)}\left(\frac{t}{t-1}\right)^3.
\end{equation}
Table~\ref{tab:prefixes} gives values of $t$ for which this upper
bound is below $2$, together with every retained prime $p_3$.
For $p_2=5$, the omitted possibilities $p_3=7,11,13$ have positive
three-prime gaps $9,5,3$, respectively, and are excluded by
\eqref{eq:prefix}. There are exactly thirteen remaining triples.

\begin{table}[htbp]
\centering\small
\caption{Complete initial-segment bounds. All primes $p_3\ge t$ are excluded.}
\label{tab:prefixes}
\begin{tabular}{@{}rrrrl@{}}
\toprule
$p_2$ & $t$ & $3p_2t^3$ & $4(p_2-1)(t-1)^3$ & Retained $p_3$\\
\midrule
5 &47&1557345&1557376&17,19,23,29,31,37,41,43\\
7 &23&255507&255552&11,13,17,19\\
11&17&162129&163840&13\\
13&15&131625&131712&none\\
\bottomrule
\end{tabular}
\end{table}

Fix one retained triple, and put $c=2B-A>0$, $x=p_4$, $y=p_5$.
Expanding the equation $Axy-2B(x-1)(y-1)=3$ gives the equivalence
\begin{equation}\label{eq:factor}
 (cx-2B)(cy-2B)=2AB-3c=:N.
\end{equation}
The four-prime initial segment has negative gap, so
$2B-cx<0$. Consequently
\[
 u=cx-2B>0,\qquad v=cy-2B>u.
\]
Every possible solution thus comes from a positive divisor $u\mid N$
with $u^2<N$, by setting
\begin{equation}\label{eq:recover}
 x=\frac{u+2B}{c},\qquad
 y=\frac{N/u+2B}{c}.
\end{equation}
Conversely, any such pair of integral primes with $p_3<x<y$ gives a
solution of the original equation. This also excludes zero and negative
factor branches, while $u^2=N$ would force $x=y$.

\begin{table}[htbp]
\centering\small
\caption{The complete gap-three factor calculation. Pairs are integral and satisfy $p_3<x<y$.}
\label{tab:factors}
\begin{tabular}{@{}lrrl p{4.05cm}@{}}
\toprule
$(p_1,p_2,p_3)$ & $B$ & $c$ & $N$ (prime factorization) & All ordered pairs $(x,y)$\\
\midrule
$(3,5,17)$ & 128 & 1 & $65277=3^{2}\cdot7253$ & $(257,65533)$; $(259,22015)$; $(265,7509)$\\
$(3,5,19)$ & 144 & 3 & $82071=3^{2}\cdot11\cdot829$ & $(97,9215)$; $(107,925)$\\
$(3,5,23)$ & 176 & 7 & $121419=3^{4}\cdot1499$ & none\\
$(3,5,29)$ & 224 & 13 & $194841=3^{2}\cdot21649$ & none\\
$(3,5,31)$ & 240 & 15 & $223155=3^{4}\cdot5\cdot19\cdot29$ & none\\
$(3,5,37)$ & 288 & 21 & $319617=3^{2}\cdot17\cdot2089$ & none\\
$(3,5,41)$ & 320 & 25 & $393525=3^{3}\cdot5^{2}\cdot11\cdot53$ & none\\
$(3,5,43)$ & 336 & 27 & $433359=3^{2}\cdot179\cdot269$ & none\\
$(3,7,11)$ & 120 & 9 & $55413=3^{2}\cdot47\cdot131$ & $(27,2079)$\\
$(3,7,13)$ & 144 & 15 & $78579=3^{2}\cdot8731$ & none\\
$(3,7,17)$ & 192 & 27 & $137007=3^{2}\cdot13\cdot1171$ & none\\
$(3,7,19)$ & 216 & 33 & $172269=3^{2}\cdot19141$ & none\\
$(3,11,13)$ & 240 & 51 & $205767=3^{3}\cdot7621$ & none\\
\bottomrule
\end{tabular}
\end{table}

Table~\ref{tab:factors} gives the complete calculation. Each displayed
factorization is into primes; the largest factor is $21649$, and trial
division through $147$ suffices to verify all the primality assertions.
Expanding the factorizations yields all positive divisors, so the
integrality conditions in \eqref{eq:recover} are finite exact checks.
For transparency, Appendix~\ref{app:divisors} lists all $95$ divisors
with $u^2\le N$. Exactly six ordered integral pairs survive. They fail
primality because, in the order displayed,
\begin{align*}
 65533&=13\cdot5041, &259&=7\cdot37, &265&=5\cdot53,\\
 9215&=5\cdot1843, &925&=5\cdot185, &27&=3\cdot9.
\end{align*}
The only further integral pair is $(25,5375)$ for the triple
$(3,5,43)$, which violates $x>43$. There are therefore no five-prime
solutions of the gap-three equation.
\end{proof}

\begin{proof}[Completion of the proof of Theorem~\ref{thm:main}]
If all five $p_i^{a_i+1}$ were at least $2R$, then
\eqref{eq:gap} would force $E=1$ or $E=3$. The congruence excludes
$E=1$ and Lemma~\ref{lem:radical} excludes $E=3$.
This contradiction proves the asserted strict inequality.
\end{proof}

\section{Verification and scope}

The arithmetic in Table~\ref{tab:factors} uses exact integer operations.
A finite certificate needs both exhaustive coverage and verification of
its individual witnesses. For a methodological comparison, the finite
geometric certificates of Zeng, Liu, Ratnavelu and Ong Seng Huat
\cite{ZengEtAl2026} (Section~3) separately reconstruct the enumerated
configurations and check the stored colorings. Here the initial-segment
bounds and divisor lists establish coverage, while congruences and
explicit proper factors reject each candidate.
A verifier can check each factorization by multiplication and trial
division, expand its divisors, and test whether both $u+2B$ and
$N/u+2B$ are divisible by $c$. The proof of completeness is supplied
by \eqref{eq:prefix}--\eqref{eq:recover}, rather than by a chosen
upper bound on the search.

A separate calculation checks the table without factoring $N$.
Solving for the last prime gives
\begin{equation}\label{eq:direct}
 y=\frac{2B(x-1)-3}{cx-2B}.
\end{equation}
Since the denominator is positive and $y>x$, it follows that
$cx^2<4Bx-2B-3<4Bx$, hence $x<4B/c$.
Testing \eqref{eq:direct} for positive denominator, integral $y$ and
$y>x$ at every integer in $p_3<x<4B/c$ involves $937$ possible
fourth values over the thirteen triples and gives the same six
ordered integral pairs. As a control outside the five-prime scope,
$3\cdot7\cdot13\cdot19=5187$ has totient $2592$ and gap $3$.
Thus Lemma~\ref{lem:radical} would be false with four in place of five.

The theorem includes $\sigma(n)=2n$ and places no upper bound on any
$a_i$, $p_i$, or $n$. The argument does not settle the conjecture for
six or more distinct prime divisors, nor a bound with $R$ in place of
$2R$. No covering proof of the five-prime assertion was identified in
the primary literature inspected through 12 September 2026; this is a
bounded literature assessment, not a claim of exhaustive historical
coverage or priority.

\paragraph{Computational and AI assistance.}
OpenAI Codex assisted with proof development, exact-integer program
construction, source comparison and manuscript preparation.
An additional AI-assisted verification pass reconstructed the finite
enumeration using a separate program. The calculations use Python~3
integer arithmetic, with rational arithmetic for auxiliary identity
checks, and involve no floating-point tests or random sampling.
The checks are not a formal proof-assistant verification or journal
peer review. The mathematical argument and its complete finite
certificate are given above and in the appendix.

\appendix
\section{Complete divisor lists}\label{app:divisors}

For each row of Table~\ref{tab:factors}, the following table lists
every positive divisor $u\mid N$ with $u^2\le N$.
An entry yields an integral pair precisely when
$u\equiv N/u\equiv-2B\pmod c$.
These lists can be reconstructed by independently ranging each prime
exponent from zero to its multiplicity in the factorization of $N$.
They contain $95$ entries in total.

\small
\begin{longtable}[c]{@{}l p{11.4cm}@{}}
\toprule
Initial triple & All small divisors $u$\\
\midrule
\endhead
$(3,5,17)$ & 1, 3, 9\\[3pt]
$(3,5,19)$ & 1, 3, 9, 11, 33, 99\\[3pt]
$(3,5,23)$ & 1, 3, 9, 27, 81\\[3pt]
$(3,5,29)$ & 1, 3, 9\\[3pt]
$(3,5,31)$ & 1, 3, 5, 9, 15, 19, 27, 29, 45, 57, 81, 87, 95, 135, 145, 171, 261, 285, 405, 435\\[3pt]
$(3,5,37)$ & 1, 3, 9, 17, 51, 153\\[3pt]
$(3,5,41)$ & 1, 3, 5, 9, 11, 15, 25, 27, 33, 45, 53, 55, 75, 99, 135, 159, 165, 225, 265, 275, 297, 477, 495, 583\\[3pt]
$(3,5,43)$ & 1, 3, 9, 179, 269, 537\\[3pt]
$(3,7,11)$ & 1, 3, 9, 47, 131, 141\\[3pt]
$(3,7,13)$ & 1, 3, 9\\[3pt]
$(3,7,17)$ & 1, 3, 9, 13, 39, 117\\[3pt]
$(3,7,19)$ & 1, 3, 9\\[3pt]
$(3,11,13)$ & 1, 3, 9, 27\\[3pt]
\bottomrule
\end{longtable}
\normalsize

\bibliographystyle{plain}
\bibliography{references}
\end{document}
